\documentclass[12pt,a4paper]{article}
\usepackage[utf8]{inputenc}
\usepackage[T1]{fontenc}
\usepackage[spanish,es-nodecimaldot]{babel}
\usepackage[a4paper,margin=2.5cm]{geometry}
\usepackage{lmodern}
\usepackage{microtype}
\usepackage{xcolor}
\usepackage{enumitem}
\usepackage{csquotes}
\usepackage[hidelinks]{hyperref}
\usepackage[backend=bibtex,style=ieee,sorting=none,maxnames=99,url=false,doi=true,isbn=false]{biblatex}
\addbibresource{fuentes.bib}
\definecolor{azul}{RGB}{35,55,80}
\setlength{\parindent}{0pt}
\setlength{\parskip}{7pt}
\hypersetup{colorlinks=true,urlcolor=azul,linkcolor=azul,citecolor=azul}

\newcommand{\ficha}[7]{%
  \subsection*{#1}
  {\color{azul}\rule{\linewidth}{0.8pt}}

  \textbf{Referencia IEEE:}

  \fullcite{#2}

  \textbf{DOI verificado:}

  \href{https://doi.org/#3}{#3}

  \textbf{Resumen del art\'iculo:}

  #4

  \textbf{Dato concreto:}

  #5

  \textbf{Nota personal:}

  #6

  \textbf{Uso en TA-IND-01:}

  #7

  \vspace{6mm}
}

\begin{document}

\begin{titlepage}
\begin{center}
{\Large\bfseries UNIVERSIDAD T\'ECNICA ESTATAL DE QUEVEDO}

\vspace{1.4cm}
{\large FACULTAD DE CIENCIAS DE LA INGENIER\'IA}

\vspace{1.4cm}
{\Large\bfseries FICHAS BIBLIOGR\'AFICAS}

\vspace{0.5cm}
{\large Fundamentos de sistemas distribuidos y arquitecturas distribuidas}

\vspace{1.4cm}
\textbf{Actividad:}

FORM-FB-01

\vspace{0.45cm}
\textbf{Asignatura:}

Aplicaciones Distribuidas

\vspace{0.45cm}
\textbf{Estudiante:}

Steven Fernando Pozo Castillo

\vspace{0.45cm}
\textbf{Docente:}

Ing. Guerrero Ulloa Gleiston Ciceron

\vspace{0.45cm}
\textbf{Curso:}

S\'eptimo Software B

\vfill
Quevedo - Ecuador

Noviembre de 2026
\end{center}
\end{titlepage}

\section{Introducci\'on}

Esta actividad presenta cinco fichas bibliogr\'aficas sobre sistemas distribuidos y arquitecturas de microservicios. Las fuentes seleccionadas estudian la migraci\'on desde sistemas monol\'iticos, el uso de inteligencia artificial durante el ciclo DevOps, los principios de dise\~no, y los patrones de despliegue y comunicaci\'on entre servicios.

El prop\'osito es reunir evidencia acad\'emica reciente y verificable que permita comprender tanto los beneficios como la complejidad de una arquitectura distribuida. Esta base bibliogr\'afica se utilizar\'a posteriormente en TA-IND-01 para analizar de manera cr\'itica un sistema distribuido real.

\section{Criterios de selecci\'on}

Se eligieron cinco art\'iculos de revistas acad\'emicas publicados entre 2021 y 2025. Todos poseen un DOI resoluble y pertenecen a editoriales reconocidas solicitadas en la actividad: IEEE, Springer Nature o Elsevier.

Para verificar cada fuente se consideraron los siguientes aspectos:

\begin{enumerate}[nosep]
  \item Correspondencia del DOI con la obra citada.
  \item Coincidencia de autores, a\~no y t\'itulo.
  \item Revista, volumen, n\'umero y p\'aginas o n\'umero de art\'iculo.
  \item Relaci\'on directa con sistemas distribuidos o microservicios.
  \item Existencia de resultados, definiciones o conclusiones comprobables.
\end{enumerate}

\section{Fichas bibliogr\'aficas}

\ficha
{Ficha 1. Aprendizaje autom\'atico para migrar monolitos a microservicios}
{trabelsi2025mlmigration}
{10.1109/TSE.2025.3603897}
{El art\'iculo estudia la dificultad de seleccionar t\'ecnicas de aprendizaje autom\'atico para migrar sistemas monol\'iticos a microservicios. Aplica una revisi\'on sistem\'atica, reportada con PRISMA, que integra investigaciones publicadas entre 2015 y 2024. Los autores clasifican las fases automatizadas, los artefactos de entrada, los algoritmos empleados, las evaluaciones y las limitaciones observadas. Los resultados muestran mayor atenci\'on a la identificaci\'on de servicios y al monitoreo, mientras el empaquetado de microservicios permanece casi sin explorar. Tambi\'en se repiten problemas de disponibilidad de datos, escalabilidad, complejidad, herramientas insuficientes y falta de bancos de prueba estandarizados. La taxonom\'ia resultante orienta decisiones y nuevas investigaciones.}
{La revisi\'on incluye \textbf{81 estudios primarios} publicados entre 2015 y 2024. Ubicaci\'on: resumen y Sec.~IV, pp.~2972--2995.}
{Esta fuente me ayud\'o a entender que separar un monolito no consiste solamente en dividir c\'odigo. Tambi\'en exige escoger datos, criterios de evaluaci\'on y herramientas adecuadas.}
{La utilizar\'e para analizar la estrategia de descomposici\'on del sistema real y discutir los l\'imites de automatizar esa decisi\'on.}

\ficha
{Ficha 2. Inteligencia artificial durante el ciclo de vida de microservicios}
{moreschini2025ai}
{10.1007/s00607-025-01432-z}
{Esta investigaci\'on organiza el uso de inteligencia artificial durante el ciclo DevOps de aplicaciones basadas en microservicios. Mediante un mapeo sistem\'atico, los autores relacionan dominios de IA, atributos de calidad y fases de desarrollo y operaci\'on. El an\'alisis encuentra una concentraci\'on marcada en eficiencia de rendimiento, escalado, planificaci\'on de recursos y detecci\'on de anomal\'ias; en cambio, las fases tempranas de desarrollo reciben menor atenci\'on. La s\'intesis propone temas que abarcan aprendizaje por refuerzo, AIOps, seguridad y reconstrucci\'on de servicios heredados. Concluye que numerosos trabajos siguen siendo prototipos y que faltan automatizaci\'on completa, explicabilidad y evaluaciones industriales suficientemente amplias para convertirlos en productos utilizables.}
{Se seleccionaron \textbf{269 art\'iculos}; aproximadamente \textbf{80\%} se concentra en eficiencia de rendimiento y el an\'alisis produce \textbf{16 temas}. Ubicaci\'on: Tabla~3, Secs.~4.3 y~5.}
{La fuente ampli\'o mi visi\'on porque conecta las decisiones arquitect\'onicas con la operaci\'on continua, la observabilidad y la respuesta ante fallos.}
{La usar\'e para evaluar el monitoreo, el escalado y la detecci\'on de fallos del sistema analizado, adem\'as de comprobar si su automatizaci\'on es explicable.}

\ficha
{Ficha 3. Migraci\'on sist\'emica y t\'ecnica hacia microservicios}
{ayas2023migration}
{10.1007/s10664-023-10308-9}
{El estudio busca explicar la migraci\'on a microservicios como un cambio prolongado que combina transformaciones organizativas y t\'ecnicas. Emplea teor\'ia fundamentada sobre dos conjuntos de evidencia: entrevistas con profesionales y discusiones de Stack Overflow. A partir de la codificaci\'on inductiva, reconstruye iteraciones de cambio, actividades recurrentes y resultados t\'ecnicos concretos. Los hallazgos indican que la arquitectura evoluciona de manera iterativa y requiere simult\'aneamente perspectivas de corto y largo plazo, adem\'as de conocimiento del negocio. Gran parte del trabajo no es separar funciones, sino crear artefactos de soporte, comunicaci\'on, datos, despliegue y operaci\'on. La propuesta muestra por qu\'e estas migraciones son sist\'emicas y multifac\'eticas.}
{El an\'alisis usa \textbf{19 entrevistas} y \textbf{215 discusiones} de Stack Overflow, y deriva \textbf{2 modos de cambio, 14 actividades y 53 resultados de soluci\'on}. Ubicaci\'on: resumen y Sec.~4.}
{Me permiti\'o distinguir entre cambiar la estructura del software y transformar la forma de trabajar de los equipos responsables del sistema.}
{La citar\'e al describir la evoluci\'on del sistema, sus decisiones de comunicaci\'on y despliegue, y el impacto de la arquitectura en los procesos de desarrollo.}

\ficha
{Ficha 4. Principios, patrones y retos de migraci\'on}
{velepucha2023survey}
{10.1109/ACCESS.2023.3305687}
{El art\'iculo examina los fundamentos que orientan una arquitectura de microservicios y el paso desde aplicaciones monol\'iticas. Mediante una revisi\'on de literatura, compara modularizaci\'on, ocultamiento de informaci\'on, principios orientados a objetos y criterios de descomposici\'on asociados con Parnas, Fowler y Newman. Adem\'as, organiza patrones arquitect\'onicos y resume ventajas y desventajas observadas en ambos estilos. Su aporte principal es mostrar que adoptar microservicios no equivale a fragmentar una aplicaci\'on arbitrariamente: cada servicio debe responder a l\'imites coherentes, capacidades del negocio y responsabilidades claras. La revisi\'on tambi\'en advierte sobre complejidad operativa, comunicaci\'on remota y gesti\'on distribuida, aspectos que pueden neutralizar los beneficios esperados.}
{La Tabla~4 indica que los microservicios favorecen el despliegue independiente, la escalabilidad y el aislamiento de fallos, pero aumentan la complejidad de comunicaci\'on, pruebas, monitoreo y consistencia de datos. Ubicaci\'on: Sec.~V, pp.~88352--88354.}
{Esta fuente me ofreci\'o un marco para no confundir modularidad con una cantidad elevada de servicios o con una fragmentaci\'on arbitraria del sistema.}
{La utilizar\'e para identificar patrones del caso real y juzgar si sus fronteras corresponden a capacidades del negocio o generan acoplamiento innecesario.}

\ficha
{Ficha 5. Patrones de despliegue y comunicaci\'on}
{aksakalli2021patterns}
{10.1016/j.jss.2021.111014}
{La revisi\'on analiza dos problemas inevitables cuando los microservicios se ejecutan como sistema distribuido: c\'omo desplegarlos y c\'omo intercambian informaci\'on. Sigue un protocolo sistem\'atico de b\'usqueda y selecci\'on, y despu\'es clasifica los enfoques informados por los estudios primarios. Los resultados distinguen alternativas de despliegue, patrones de comunicaci\'on y retos asociados con descubrimiento, balanceo de carga, replicaci\'on de datos y llamadas remotas. Publicaci\'on/suscripci\'on aparece como el mecanismo de comunicaci\'on m\'as frecuente, mientras una instancia de servicio por contenedor domina el despliegue. Los autores concluyen que no existe una alternativa universal: la selecci\'on depende del contexto, de los atributos de calidad buscados y de la capacidad operativa disponible.}
{De \textbf{239 trabajos revisados}, \textbf{38} fueron estudios primarios; se identificaron \textbf{3 enfoques de despliegue, 7 patrones de comunicaci\'on, 8 retos de despliegue y 6 de comunicaci\'on}. Ubicaci\'on: resumen, Sec.~4 y conclusiones.}
{Esta lectura conect\'o la estructura l\'ogica de los servicios con la infraestructura donde realmente se ejecutan y se comunican.}
{La emplear\'e para describir c\'omo se comunican y despliegan los servicios, y evaluar si sus decisiones favorecen disponibilidad y escalabilidad sin complejidad excesiva.}

\newpage
\printbibliography[title={Referencias bibliogr\'aficas}]

\section*{Declaraci\'on de uso de inteligencia artificial}

Utilic\'e ChatGPT (OpenAI) para localizar fuentes candidatas, comprobar metadatos, organizar el documento en LaTeX y apoyar la redacci\'on inicial. Cada referencia se verific\'o abriendo su DOI y contrastando en la publicaci\'on original el t\'itulo, los autores, el a\~no, la revista, el volumen, el n\'umero y las p\'aginas o n\'umero de art\'iculo. Los datos concretos se contrastaron con el resumen, las tablas o las secciones indicadas. La revisi\'on final, la interpretaci\'on de las fuentes y la responsabilidad por el texto presentado corresponden al estudiante.

\end{document}
